% ==============================================================================
% Section 6: Discussion
% ==============================================================================

\section{Discussion}
\label{sec:discussion}

\subsection{Recovery without intervention}

Across all four forecast scenarios, the model predicts that
\textit{Pycnopodia helianthoides} populations will stabilize at
approximately 94--95\% below pre-disease levels through mid-century
(F01: 94.6\% decline; F04: 94.2\% decline).  No combination of climate
trajectory and evolutionary capacity produces meaningful coast-wide
recovery to pre-disease abundance.  The absolute population floor shifts
by tens of thousands of individuals between scenarios, but the
qualitative conclusion is unchanged: the species' new normal is a
small fraction of its former abundance.

The mechanism underlying this result is what we term the
\textbf{seasonal disease ratchet}.  During winter, falling water
temperatures push \textit{V.\ pectenicida} into VBNC dormancy,
disease pressure abates, and surviving adults reproduce.  Larval
recruitment partially replenishes depleted sites.  But each summer,
rising temperatures reactivate the pathogen, the environmental pool
rebuilds through the shedding--infection positive feedback, and newly
recruited cohorts are culled alongside the survivors.  The net effect
is a population that oscillates around a low equilibrium, gaining
recruits in winter and losing them in summer, unable to escape the
ratchet.  This pattern is especially pronounced from Washington
southward, where summer SST consistently exceeds the VBNC activation
threshold and disease activity spans 6--10 months of the year.

\subsection{The Oregon oasis paradox}

Oregon presents a complex case study in episodic recovery dynamics. 
While 2050 endpoint snapshots suggest 17.6\% of pre-disease capacity 
is retained under F01, the 5-year average reveals only 5.4\% with 121\% 
coefficient of variation---indicating boom-bust cycles rather than 
sustained recovery. Alaska's Fairweather-North region shows more 
reliable persistence (6.8\% endpoint, 9.4\% 5-year average). 
Nevertheless, Oregon's episodic peaks demonstrate what is possible 
under favorable conditions.  Three
features converge to favor Oregon:

\begin{enumerate}
  \item \textbf{Temperature regime.}  Oregon's SST occupies a Goldilocks
    zone---warm enough that the growing and reproductive season is
    adequate for population maintenance, but cool enough that the
    pathogen enters dormancy for several months each winter.  The final
    pathogen thermal threshold at Oregon sites ($T_\mathrm{vbnc} \approx
    10.7$\textdegree C) reflects moderate cold-adaptation, but winter
    SST still drops well below this value, providing genuine disease
    respite.

  \item \textbf{Open-coast flushing.}  Oregon's exposed outer coast
    scores low on the enclosedness metric, yielding high flushing rates
    that continuously dilute the environmental pathogen pool.  Unlike
    the semi-enclosed waters of Puget Sound or British Columbia's fjords,
    where pathogen concentrations build in poorly flushed basins, Oregon's
    sites experience continuous hydrodynamic renewal.

  \item \textbf{Moderate virulence.}  The local pathogen community has
    evolved to intermediate virulence ($v_\mathrm{local} \approx 0.14$),
    low enough that hosts survive long enough to reproduce between
    outbreaks, but high enough that the disease exerts meaningful selection
    for resistance.  Oregon populations have evolved resistance to
    $r \approx 0.38$ by 2050---high enough to meaningfully slow infection
    acquisition.
\end{enumerate}

Under high-end warming (SSP5-8.5, F04), Oregon's episodic booms diminish  
(endpoint drops to 10.4\%) because warmer summers extend the pathogen
activity window.  This underscores that Oregon's oasis dynamics are
thermally mediated and climate-contingent.  For conservation managers,
Oregon represents a double-edged opportunity: its episodic booms 
demonstrate recovery potential and provide valuable monitoring targets, 
but the underlying boom-bust volatility means Oregon cannot be relied 
upon for stable long-term persistence. Alaska's Fairweather regions 
may offer more reliable reintroduction prospects.

\subsection{Alaska's conditional optimism}

Alaska presents a more complicated picture.  The model predicts only
2--10\% recovery for Alaskan regions across climate scenarios, but
we know this is an underestimate: field surveys document substantially
higher recovery in Prince William Sound (estimated ${\sim}$50\%;
\citep{gehman2025}).  The model's Alaskan underprediction is its
single largest calibration shortcoming (Section~\ref{sec:calibration})
and likely reflects processes not fully captured in the current
framework---including freshwater surface lenses that suppress
\textit{Vibrio} in fjord environments, localized upwelling, or
higher background cold-tolerance in northern host genotypes.

Setting aside the absolute numbers, the model's \textit{directional}
prediction for Alaska is noteworthy.  Under SSP5-8.5, Fairweather
North recovery increases from 6.8\% (F01) to 9.9\% (F04)---a 45\%
relative improvement.  Warming extends the growing season and
reproductive window in cold-water Alaskan habitats without
proportionally extending the disease season, because SST in these
regions remains near or below the VBNC floor for much of the year
even under high emissions.  If real Alaskan populations are indeed
doing substantially better than the model shows, then climate
change---paradoxically---may create further opportunity for northern
populations by expanding favorable thermal habitat.

This has a practical implication: Alaskan populations may be the
species' best long-term reservoir, and monitoring investments in
Prince William Sound and the Fairweather coast are well justified.

\subsection{Southern California as a conservation priority}

At the opposite end of the range, the model delivers an unambiguous
verdict.  Southern California (CA-S) and Baja California retain
effectively zero population across all scenarios (CA-S: 0.01\% recovery,
${\sim}$52 individuals coast-wide; Baja: functionally extinct with zero
individuals).  Central California (CA-C) retains ${\sim}$12\% under
moderate warming but declines to 7.6\% under SSP5-8.5, and this trend
is downward through 2050 with no indication of stabilization.

Year-round pathogen activity, high evolved virulence
($v_\mathrm{local} \approx 0.30$), and the absence of any seasonal disease
respite make natural recovery effectively impossible in these regions.
The environmental pathogen pool never decays to levels low enough for
naïve recruits to survive their first summer.  Even with strong
selection (resistance evolves to $r \approx 0.43$--0.45 in CA-C by 2050),
the host cannot evolve fast enough to escape the pathogen's demographic
trap.

These findings motivate the upcoming F05--F08 reintroduction scenarios,
which will test whether captive-bred juveniles with elevated resistance
can establish self-sustaining populations in southern waters.  Our
results suggest that any reintroduction effort south of Point
Conception will require \textit{sustained} releases over multiple
years, because single-cohort releases will likely be overwhelmed by the
ambient disease pressure before the population can reach densities
sufficient to buffer annual losses.

\subsection{The evolutionary dead end}

One of the most striking---and sobering---results is that natural
selection, while clearly operating, is insufficient to rescue
\textit{Pycnopodia} from SSWD.  Host resistance evolves substantially
in hard-hit populations: central California's mean resistance increases
from $r = 0.15$ to $r = 0.43$ over 38 simulated years, nearly tripling
the individual-level hazard reduction.  Yet $r = 0.43$ still means that
58\% of the pathogen force of infection penetrates the host's defenses.
In a high-pathogen environment, this translates to effective daily
infection probabilities that remain well above 20\%.  Resistance buys
time, not survival.

The fundamental asymmetry is temporal.  \textit{V.\ pectenicida}
can adapt its thermal tolerance and virulence on daily-to-seasonal
timescales through community-level selection among trillions of
bacteria.  The host evolves its defense traits once per generation,
with a minimum generation time of several years and effective
population sizes ($N_e$) reduced far below census counts by
sweepstakes reproduction \citep{hedgecock2011}.  This mismatch means
the pathogen can track environmental changes orders of magnitude faster
than the host can respond.

Compounding this asymmetry, the populations under strongest selection
are also losing the genetic raw material needed for future adaptation.
Additive genetic variance for resistance in central California drops
88\% (from $V_A = 0.026$ to $V_A = 0.003$) as population bottlenecks
eliminate rare alleles.  Southern populations are being driven into an
evolutionary corner: strong selection, but nothing left to select on.

The implication for conservation is clear.  Captive breeding programs
that screen for high-resistance broodstock \citep{hodin2021} and
maintain genetic diversity through managed crosses are not merely
helpful supplements to natural recovery---they may be the \textit{only}
path to re-establishing \textit{Pycnopodia} in the southern half of
its range, where natural selection alone leads to an evolutionary dead
end.

\subsection{Limitations}

We highlight several important caveats:

\begin{itemize}
  \item \textbf{Alaska underprediction.}  As discussed, the model
    produces 2--10\% Alaskan recovery against ${\sim}$50\% field
    estimates.  This is the model's main weakness and likely reflects
    missing local processes (fjord stratification, freshwater lensing)
    that suppress \textit{Vibrio} activity in these environments.

  \item \textbf{Genetic architecture.}  The 51-locus, three-trait
    framework is borrowed from GWAS data on \textit{Pisaster ochraceus}
    \citep{schiebelhut2018}.  The true genetic basis of SSWD resistance
    in \textit{Pycnopodia} is unknown and may involve different numbers
    of loci, epistatic interactions, or trait architectures not captured
    by the additive model.

  \item \textbf{Climate projections.}  The delta method applied to 11
    CMIP6 reference sites cannot capture fine-scale oceanographic
    features---local upwelling intensification, mesoscale eddies, or
    freshwater-driven thermal anomalies---that may meaningfully alter
    SST trajectories at specific sites.

  \item \textbf{Simplified dispersal.}  Larval and pathogen dispersal
    use distance-decay kernels without explicit ocean current data.
    Real connectivity patterns are shaped by coastal currents, seasonal
    reversals, and mesoscale features that our kernels cannot resolve.

  \item \textbf{Single pathogen.}  We model \textit{V.\ pectenicida}
    as the sole disease agent.  The true SSWD etiology may involve
    co-infections, opportunistic secondary pathogens, or microbiome
    dysbiosis that amplifies the primary infection.

  \item \textbf{Resistance as a scalar trait.}  Real disease resistance
    arises from complex immune networks with epistasis, dominance, and
    genotype-by-environment interactions.  Our additive scalar
    approximation captures the broad dynamics of selection but not the
    full complexity of host--pathogen coevolution at the molecular level.
\end{itemize}

\subsection{What comes next}

Three immediate extensions of this work are planned:

\begin{description}
  \item[Reintroduction scenarios (F05--F08):]  We will simulate captive
    breeding releases at depleted southern sites, testing different
    release sizes, frequencies, broodstock resistance levels, and site
    selection strategies.  These scenarios will directly inform the
    operational decisions facing the \textit{Pycnopodia} recovery
    program.

  \item[Sensitivity analysis:]  A formal Morris screening followed by
    Sobol variance-based decomposition will identify which model
    parameters most strongly influence population outcomes.  This will
    focus future calibration effort where it matters most and quantify
    the uncertainty envelope around our forecasts.

  \item[Publication manuscript:]  The present report will be condensed
    and submitted for peer review, providing the broader marine
    conservation community with a quantitative framework for
    \textit{Pycnopodia} recovery planning.
\end{description}

Taken together, the forecast results paint a picture that is challenging
but not hopeless.  Natural recovery alone will not restore
\textit{Pycnopodia} to its former abundance anywhere in its range, and
southern populations face functional extinction without intervention.
But the model also identifies genuine bright spots---Oregon's thermal
refuge, Alaska's cold-water reservoir, and the steady (if slow)
accumulation of resistance alleles in surviving populations---that
provide conservation managers with concrete targets for monitoring,
protection, and eventual reintroduction.
